
\subsection{Lower endpoint and continuity}

This subsection proves the \(r\downarrow0\) assertion of
Proposition~\ref{prop:phase-curve-properties}. A single large positive Poisson
point governs the one-sided crossing probability.
\begin{lemma}
\label{lem:endpoint-one-large-jump}
For one copy of the limiting Poisson field,
\[
    R_{1-\eps}
    :=
    \mathbb P\Bigl\{
        \sup_{1-\eps\le\theta<1}S(\theta)>0
    \Bigr\}
    \sim \eps
    \qquad(\eps\downarrow0).
\]
\end{lemma}

\begin{proof}
Write $h=1-\theta$. Split $S(1-h)=A_h+B_h$, where
\[
    A_h
    =
    \int_0^\infty e^{(1-h)x}\{N-\mu\}(dx),
    \qquad
    B_h
    =
    \int_{-\infty}^0 e^{(1-h)x}\{N-\mu\}(dx).
\]
Enumerate the points of \(N\cap[0,\infty)\) as \(X_j\), and put
\(W_j=e^{X_j}\). Then the \(W_j\) form a Poisson point process on
\([1,\infty)\) with intensity \(w^{-2}\,dw\), and hence
\[
    A_h=\sum_j W_j^{1-h}-\frac1h .
\]
The negative half-line is negligible at the scale \(1/h\). Indeed, for every
fixed \(q\ge2\), the compensated Poisson integrals defining \(B_h\) and
\(\partial_hB_h\) have uniformly bounded \(q\)th moments for
\(0\le h\le1/4\). The fundamental theorem of calculus therefore gives
\(\mathbb E\sup_{0\le h\le1/4}|B_h|^q<\infty\).
Consequently, with \(K_\eps=\eps^{-1/3}\) and any \(q>3\),
\[
    \mathbb P\Bigl\{
        \sup_{0<h\le\eps}|B_h|>K_\eps
    \Bigr\}
    =
    o(\eps).
\]
It therefore suffices to estimate the positive-half-line crossing probability
with a perturbation of size \(K_\eps\). For each fixed sign
$\sigma\in\{-1,1\}$, consider
\[
    \mathbb P\Bigl\{
        \sup_{0<h\le\eps}
        \Bigl(\sum_j W_j^{1-h}-\frac1h\Bigr)
        >\sigma K_\eps
    \Bigr\}.
\]
The sign does not affect the first-order answer, since $K_\eps=o(\eps^{-1})$.

For fixed $h$, put $Y_j=W_j^{1-h}$. Then the $Y_j$ form a Poisson cloud on $[1,\infty)$ with tail intensity
\[
    \nu_h(y,\infty)=y^{-1/(1-h)}.
\]
This regularly varying tail is subexponential; we need its one-large-jump
estimate uniformly over \(0<h\le\eps\)
\cite[Chapter~4]{foss2013introduction}. Write
\[
 \beta=(1-h)^{-1},\qquad
 x=x_h=h^{-1}+\sigma K_\eps,\qquad
 \delta=1/\log x.
\]
The total number
of points is \(\operatorname{Pois}(1)\), and
\[
    \nu_h(y,\infty)=y^{-\beta},
    \qquad
    \int_1^{\delta x}y\,\nu_h(dy)
    =
    \beta\int_1^{\delta x}y^{-\beta}\,dy
    \le C\log x
\]
uniformly for \(0<h\le\eps\), once \(\eps<1/3\). If
\(\sum_jY_j>x\) but \(\max_jY_j\le x\), then at least one of the following
events occurs:
\begin{enumerate}
\item there is a point in \(((1-\delta)x,x]\);
\item there are at least two points above \(\delta x\);
\item every point is at most \(\delta x\), and the total number of points
      exceeds \(\delta^{-1}\);
\item there is exactly one point in \((\delta x,(1-\delta)x]\), and the sum
      of the points below \(\delta x\) exceeds \(\delta x\).
\end{enumerate}
The Poisson formula, followed by Markov's inequality in the last case, bounds
the four probabilities by
\[
    C\delta x^{-\beta},\qquad
    C(\delta x)^{-2\beta},\qquad
    \mathbb P\{\operatorname{Pois}(1)>\delta^{-1}\},\qquad
    C(\delta x)^{-\beta}\frac{\log x}{\delta x},
\]
respectively. Since \(1/2\le hx\le2\), their sum is \(a_\eps h\) for a
deterministic modulus \(a_\eps\downarrow0\), uniformly over
\(0<h\le\eps\). Thus
\[
    \mathbb P\left\{
        \sum_jY_j>x_h,\
        \max_jY_j\le x_h
    \right\}
    \le a_\eps h.
\]

To pass from fixed \(h\) to the supremum, choose
\(\delta_\eps\downarrow0\) so slowly that
\(a_\eps/\delta_\eps\to0\) and \(K_\eps\eps=o(\delta_\eps)\), and use the mesh
\(h_k=\eps(1+\delta_\eps)^{-k}\). Fix a large numerical constant \(C\) and
define the enlarged one-large-jump event
\[
 \mathcal J_\eps^{\rm enl}
 =
 \left\{\max_jW_j>
 \inf_{0<h\le\eps}
 \bigl[(1-C\delta_\eps)x_h\bigr]^{1/(1-h)}
 \right\}.
\]
The infimum is attained at \(h=\eps\) for all small \(\eps\). Indeed, the
derivative of
\(\log\{[(1-C\delta_\eps)x_h]^{1/(1-h)}\}\) has the sign of
\[
 \log\{(1-C\delta_\eps)x_h\}
 -\frac{1-h}{h(1+\sigma K_\eps h)},
\]
which is negative uniformly on \(0<h\le\eps\), because
\(K_\eps\eps=o(\delta_\eps)\) and
\(h^{-1}\gg\log(1/h)\). The intensity of the enlarged event is therefore
\[
 \bigl[(1-C\delta_\eps)x_\eps\bigr]^{-1/(1-\eps)}
 =(1+O(\delta_\eps)+o(1))\eps.
\]
Thus the enlarged event differs from the exact one-large-jump event by probability
\(o(\eps)\).

On the complement of \(\mathcal J_\eps^{\rm enl}\), every atom lies below
the enlarged threshold. If \(h_{k+1}<h\le h_k\), then
\(W^{1-h_{k+1}}\ge W^{1-h}\) for every \(W\ge1\), while
\(x_{h_{k+1}}/x_h=1+O(\delta_\eps)\) uniformly in \(k\). Thus a crossing
somewhere in the \(k\)th cell gives, at \(h_{k+1}\), a subcritical-sum event
at threshold \((1-C\delta_\eps)x_{h_{k+1}}\), after increasing \(C\) if
necessary. The
fixed-\(h\) calculation above applies uniformly to this threshold as well,
because \(h_{k+1}(1-C\delta_\eps)x_{h_{k+1}}\) stays between two fixed positive
constants; enlarge the modulus \(a_\eps\downarrow0\) if necessary to cover
both thresholds. On the complement of \(\mathcal J_\eps^{\rm enl}\), its largest
atom is below this same threshold. The probability of the mesh event is
therefore at most \(Ca_\eps h_k\). Since
\(\sum_kh_k=O(\eps/\delta_\eps)\), the total contribution is \(o(\eps)\).

Thus, to first order, a crossing is caused by one atom. For either sign
\(\sigma\), differentiation shows
that the required atom threshold is decreasing on \(0<h\le\eps\), for all
small \(\eps\), and is therefore minimized at \(h=\eps\):
\[
    \inf_{0<h\le\eps}
    (h^{-1}+\sigma K_\eps)^{1/(1-h)}
    =
    (\eps^{-1}+\sigma K_\eps)^{1/(1-\eps)}
    =
    \eps^{-1/(1-\eps)}(1+o(1)).
\]
The intensity of atoms above this threshold is
\[
    \int_{(\eps^{-1}+\sigma K_\eps)^{1/(1-\eps)}}^\infty
    w^{-2}\,dw
    =
    (\eps^{-1}+\sigma K_\eps)^{-1/(1-\eps)}
    \sim\eps.
\]
Since this intensity tends to zero, the probability of at least one such atom is also asymptotic to $\eps$. Combining this with the $o(\eps)$ contribution of the negative half-line and of the subcritical positive cloud proves $R_{1-\eps}\sim\eps$.
\end{proof}

Taking \(\eps=\sqrt r\) in Lemma~\ref{lem:endpoint-one-large-jump} and using
the independence of \(S_+\) and \(S_-\) gives
\[
    p(r)=\bigl(1-R_{1-\sqrt r}\bigr)^2,
    \qquad
    1-p(r)\sim2\sqrt r
    \quad(r\downarrow0).
\]
In particular, \(p(r)\to1=p(0)\) at the lower endpoint.

\subsection{Upper-endpoint persistence exponent}

This subsection proves the \(r\uparrow1/4\) assertion of
Proposition~\ref{prop:phase-curve-properties}. Lemma
~\ref{lem:poisson-left-endpoint} gives convergence to zero but not its rate;
the sharper estimate comes from the persistence exponent of a stationary
Gaussian process with sech covariance.

\begin{lemma}
\label{lem:sech-persistence}
Let $G=(G(t))_{t\in\R}$ be the centered stationary Gaussian process with covariance
\[
    \E G(s)G(t)
    =
    \operatorname{sech}\!\left(\frac{t-s}{2}\right).
\]
For every deterministic $a_T\to0$,
\begin{equation}
\label{eq:sech-persistence}
    \mathbb P\{G(t)\le a_T\text{ for all }0\le t\le T\}
    =
    \exp\!\left\{-\frac{3}{16}T+o(T)\right\}.
\end{equation}
\end{lemma}

\begin{proof}
To keep track of the normalization, we derive the stated exponent directly
from the cited results. Dembo--Poonen--Shao--Zeitouni introduce, in
\cite[equation~(1.3)]{DemboPoonenShaoZeitouni2002}, the constant
\[
    b
    =
    -4\lim_{T\to\infty}
    \frac1T
    \log
    \mathbb P\{G(t)\le0\text{ for all }0\le t\le T\}.
\]
Here their process \(Y\) is exactly the centered stationary Gaussian process
with correlation \(\operatorname{sech}(t/2)\); see
\cite[equation~(1.4) and the sentence following it]{DemboPoonenShaoZeitouni2002}.
Their Theorem~1.1 is the no-real-zero statement, and Theorems~1.3 and 1.4
identify the same constant \(b\) as the persistence exponent obtained after
the standard reduction from the polynomial near \(\pm1\) to the stationary sech
process.

FitzGerald--Tribe--Zaboronski rigorously compute this constant. In \cite[Proposition~5]{FitzGeraldTribeZaboronski2022} they prove that the Kac-polynomial persistence probability satisfies
\[
    \lim_{n\to\infty}
    \frac{\log p_{2n+1}}{\log n}
    =
    -\frac34.
\]
Their proof explicitly invokes the Dembo--Poonen--Shao--Zeitouni characterization: see \cite[equations~(19)--(20)]{FitzGeraldTribeZaboronski2022}. Their Pfaffian formula for the sech-kernel gap probability then gives \cite[equation~(21)]{FitzGeraldTribeZaboronski2022}
\[
    \log \mathbb P\{G(t)\ne0\text{ for all }0\le t\le T\}
    =
    -\frac{3}{16}T+o(T).
    \tag{*}
\]
Because $G$ is symmetric and its zeros are almost surely simple, the no-zero event on $[0,T]$ is the disjoint union of the positive- and negative-persistence events, which have equal probability. Simplicity follows, for example, from Bulinskaya's lemma, since the sech process is smooth and $(G(t),G'(t))$ is nondegenerate for each fixed $t$. Multiplication by the constant factor $2$ does not affect the exponential rate, so $(*)$ gives
\[
    \mathbb P\{G(t)\le0\text{ for all }0\le t\le T\}
    =
    \exp\!\left\{-\frac{3}{16}T+o(T)\right\}.
\]
The same value was identified earlier in the physics literature by Poplavskyi--Schehr \cite{PoplavskyiSchehr2018}.

We also need persistence at a vanishing deterministic level
\(a_T=o(1)\), rather than exactly at zero.
Feldheim--Feldheim--Mukherjee prove existence of fixed-level persistence
exponents in \cite[Theorem~1]{FeldheimFeldheimMukherjee2025} and continuity of
the exponent in the level in \cite[Theorem~4(II)]{FeldheimFeldheimMukherjee2025}.
The sech process is covered by their hypotheses: its spectral density is
proportional to \(\operatorname{sech}(\pi\lambda)\), hence is absolutely
continuous, positive at the origin, and has exponential tails. Their statements
use strict inequalities, but continuity in the level lets us squeeze the weak
event needed here between the two fixed levels \(\pm\delta\). For every fixed
\(\delta>0\) and all large \(T\),
\[
    \{G(t)\le-\delta\text{ for all }0\le t\le T\}
    \subseteq
    \{G(t)\le a_T\text{ for all }0\le t\le T\}
    \subseteq
    \{G(t)\le\delta\text{ for all }0\le t\le T\}.
\]
Taking logarithms, dividing by $T$, and then letting $\delta\downarrow0$ gives \eqref{eq:sech-persistence}.
\end{proof}

The factor $3/16$ is the persistence exponent in logarithmic time $T$. In Proposition~\ref{prop:upper-endpoint-persistence}, $T$ will be comparable to $\log(1/\beta)$, so the one-sided no-crossing probability is $\beta^{3/16+o(1)}$; the two-sided one-atom probability is its square, hence has exponent $3/8$. This endpoint refinement is not needed for the phase diagram in Theorem~\ref{thm:exponent-phase-diagram}, but it identifies the rate at which the phase curve reaches zero.

\begin{proposition}
\label{prop:upper-endpoint-persistence}
As \(\alpha\uparrow1/4\),
\[
    p(\alpha)
    =
    \left(\frac14-\alpha\right)^{3/8+o(1)}
    .
\]
Equivalently, if
\[
    q(\alpha)
    :=
    \mathbb P\left\{
        \sup_{\theta_\alpha\le\theta<1}S(\theta)\le0
    \right\},
    \qquad
    \theta_\alpha=1-\sqrt\alpha,
\]
then $q(\alpha)=(1/4-\alpha)^{3/16+o(1)}$.
\end{proposition}

\begin{proof}
The proof has two main inputs. Near \(\theta=1/2\), the contribution of the
negative half-line has a Gaussian tangent process, and in logarithmic distance
from the endpoint this tangent process has the sech covariance of
Lemma~\ref{lem:sech-persistence}. Away from \(\theta=1/2\), the remaining
interval has only a subexponential persistence cost.

Put $\beta=\theta_\alpha-1/2=1/2-\sqrt\alpha$. Then $\beta=(1/4-\alpha)/(1/2+\sqrt\alpha)\sim1/4-\alpha$ as $\alpha\uparrow1/4$. Define the one-sided no-crossing probability
\[
    F(\beta)
    =
    \mathbb P\left\{
        S\left(\frac12+h\right)\le0
        \text{ for every }\beta\le h<\frac12
    \right\}.
\]
By definition, $q(\alpha)=F(\beta)$ and $p(\alpha)=F(\beta)^2$. It is therefore enough to prove
\begin{equation}
\label{eq:F-beta-persistence}
    F(\beta)=\beta^{3/16+o(1)}
    \qquad(\beta\downarrow0).
\end{equation}

Split the Poisson field as $S(1/2+h)=T_h+U_h$, where
\[
    T_h=\int_{-\infty}^0 e^{(1/2+h)x}\{N-\mu\}(dx),
    \qquad
    U_h=\int_0^\infty e^{(1/2+h)x}\{N-\mu\}(dx).
\]
The negative half-line is the singular part as $h\downarrow0$: its variance is
of order $1/h$. The positive half-line stays tight on the short window
$0<h\le h_0$ and will only move the persistence threshold by $o(1)$ after
normalization by $\sqrt h$.

After the change of variables $r=e^{-x}$ on $x<0$,
\[
    T_h
    =
    \int_1^\infty r^{-1/2-h}\{\Pi(dr)-dr\},
\]
where $\Pi$ is a unit-rate Poisson process on $[1,\infty)$. Let $M(r)=\Pi([1,r])-(r-1)$. Integration by parts gives
\[
    T_h
    =
    \left(\frac12+h\right)
    \int_1^\infty r^{-3/2-h}M(r)\,dr.
\]

We use the KMT approximation for a Poisson counting process. Corollary~5.3
of Chapter~7 in Ethier--Kurtz
\cite[pp.~358--359]{ethier1986markov} gives a coupling of \(M\) to a Brownian
motion \(B\) on \([1,\infty)\), with \(B(1)=0\) and covariance
\(\E B(r)B(s)=(r\wedge s)-1\), for which
\[
    \Gamma
    :=
    \sup_{r\ge1}
    \frac{|M(r)-B(r)|}{C_0(1+\log r)}
\]
has a finite exponential moment. Consequently, for constants \(C,c>0\) and
every \(y\ge1\),
\begin{equation}
\label{eq:poisson-kmt}
    \mathbb P\left\{
        \sup_{r\ge1}
        \frac{|M(r)-B(r)|}{C_0(1+\log r)}>y
    \right\}
    \le
    Ce^{-cy}.
\end{equation}
This is the Poisson-process form of the Komlós--Major--Tusnády strong
approximation \cite{komlos1975approximation,komlos1976approximation}.
For completeness, the weighted infinite-horizon form also follows from the
finite-horizon estimate by applying it at horizons \(e^{m+1}\): on the block
\(e^m\le r\le e^{m+1}\), a violation at level \(y\) requires an error of
order \(y(m+1)\), and summing the resulting
\(Ce^{-cy(m+1)}\) bounds over \(m\ge0\) gives \eqref{eq:poisson-kmt}.
Define the Gaussian comparison field
\[
    \widetilde T_h
    =
    \int_1^\infty r^{-1/2-h}\,dB(r)
    =
    \left(\frac12+h\right)
    \int_1^\infty r^{-3/2-h}B(r)\,dr.
\]
On the complement of the event in \eqref{eq:poisson-kmt},
\[
    \sqrt{2h}\,|T_h-\widetilde T_h|
    \le
    C y\sqrt h
    \int_1^\infty (1+\log r)r^{-3/2-h}\,dr
    \le
    C y\sqrt h.
\]
Let $h_0=(\log(1/\beta))^{-4}$ and $T_\beta=\log(h_0/\beta)$. Taking $y=AT_\beta$ with $A$ fixed and large gives, outside an event of probability $O(e^{-cAT_\beta})$,
\begin{equation}
\label{eq:poisson-brownian-window}
    \sup_{\beta\le h\le h_0}
    \sqrt{2h}\,|T_h-\widetilde T_h|
    \le
    \Delta_\beta,
\end{equation}
where $\Delta_\beta=C_A T_\beta\sqrt{h_0}=o(1)$.

Now parametrize $h=\beta e^t$, $0\le t\le T_\beta$, and set $G_\beta(t)=\sqrt{2\beta e^t}\,\widetilde T_{\beta e^t}$.
This process has exactly the law of $G$ restricted to $[0,T_\beta]$, since
\[
    \E G_\beta(s)G_\beta(t)
    =
    2\sqrt{\beta e^s\,\beta e^t}
    \int_1^\infty r^{-1-\beta e^s-\beta e^t}\,dr
    =
    \frac{2e^{(s+t)/2}}{e^s+e^t}
    =
    \operatorname{sech}\!\left(\frac{t-s}{2}\right).
\]
Thus the logarithmic distance from the endpoint, rather than the original
parameter $h$, is the natural time variable. The covariance computation is
exact; the only approximation in the finite window is the KMT replacement of
the Poisson counting process by Brownian motion.

We first estimate the finite window
\[
    E_\beta(h_0)
    =
    \left\{
        S\left(\frac12+h\right)\le0
        \text{ for every }\beta\le h\le h_0
    \right\}.
\]
For $0<h\le h_0<1/4$,
\[
    U_h
    =
    \sum_{x_j>0}e^{(1/2+h)x_j}
    -
    \frac{1}{1/2-h}
    \ge
    -4.
\]
Thus $E_\beta(h_0)$ implies $T_h\le4$ for every $\beta\le h\le h_0$, hence $\sqrt{2h}T_h\le4\sqrt{2h_0}$ on the same window. Combining this implication with \eqref{eq:poisson-brownian-window}, and choosing $A$ in \eqref{eq:poisson-kmt} large enough that the coupling error is negligible on the $\exp\{-(3/16)T_\beta\}$ scale, gives
\[
    \mathbb P(E_\beta(h_0))
    \le
    \mathbb P\{G(t)\le4\sqrt{2h_0}+\Delta_\beta\text{ for all }0\le t\le T_\beta\}
    +o\!\left(e^{-3T_\beta/16}\right).
\]
By \eqref{eq:sech-persistence}, this upper bound is $\exp\{-(3/16)T_\beta+o(T_\beta)\}$.

For the lower bound, let
\[
    V=\sum_{x_j>0} e^{3x_j/4}.
\]
Since $\E V=\int_0^\infty e^{-x/4}\,dx=4$, choose $K<\infty$ such that $c_K:=\mathbb P\{V\le K\}>0$. For all small $\beta$, $h_0<1/4$, so $U_h\le V$ for every $0<h\le h_0$. On $\{V\le K\}$, the event $\{T_h\le -K\text{ for all }\beta\le h\le h_0\}$ implies $E_\beta(h_0)$. The positive and negative half-line Poisson processes are independent, so
\[
    \mathbb P(E_\beta(h_0))
    \ge
    c_K\,
    \mathbb P\{T_h\le -K\text{ for all }\beta\le h\le h_0\}.
\]
On the coupling event \eqref{eq:poisson-brownian-window}, the condition
\[
    G_\beta(t)\le -K\sqrt{2h_0}-\Delta_\beta
    \qquad(0\le t\le T_\beta)
\]
implies $T_h\le -K$ throughout $\beta\le h\le h_0$. The same choice of $A$ makes the complement of the coupling event negligible on this scale. Another application of \eqref{eq:sech-persistence}, with the fixed factor $c_K$ absorbed into the $o(T_\beta)$ term, therefore gives
\[
    \mathbb P(E_\beta(h_0))
    \ge
    \exp\!\left\{-\frac{3}{16}T_\beta+o(T_\beta)\right\}.
\]
Together with the upper bound,
\begin{equation}
\label{eq:finite-window-persistence}
    \mathbb P(E_\beta(h_0))
    =
    \exp\!\left\{-\frac{3}{16}T_\beta+o(T_\beta)\right\}
    =
    \beta^{3/16+o(1)},
\end{equation}
because $T_\beta=\log(1/\beta)+o(\log(1/\beta))$.

The upper bound in \eqref{eq:F-beta-persistence} follows immediately from $F(\beta)\le\mathbb P(E_\beta(h_0))$. For the lower bound, use the Harris--FKG inequality for Poisson processes: two decreasing events of a Poisson configuration are positively correlated; see, for example, \cite[Theorem~20.4]{kallenberg2017random}. For an interval $I\subset(0,1/2)$, the event
\[
    \left\{
        S\left(\frac12+h\right)\le0
        \text{ for every }h\in I
    \right\}
\]
is decreasing in the Poisson configuration, because adding a point increases $S(\theta)$ for every $\theta$. Hence
\[
    F(\beta)
    \ge
    \mathbb P(E_\beta(h_0))\,
    \mathbb P\left\{
        S\left(\frac12+h\right)\le0
        \text{ for every }h_0\le h<\frac12
    \right\}.
\]
We now show that the second factor is subexponential on the
$\log(1/\beta)$ scale.

The interval \([h_0,1/2)\), which lies away from the singular endpoint
\(h=0\), contributes only a subexponential factor. We cover its
initial portion by a slowly expanding sequence of windows. Each window has
the same sech-persistence estimate with a new left endpoint, and the sum of
their logarithmic costs is only $O(\log(1/h_0))$.

Choose \(h_*>0\) so small that
\((\log(1/h))^{-4}>h\) for \(0<h\le h_*\). Starting from the present \(h_0\),
define \(h_{j+1}=(\log(1/h_j))^{-4}\) and stop when \(h_j\) first exceeds
\(h_*\). The finite-window lower bound just proved applies uniformly with
\(\beta\) replaced by any sufficiently small left endpoint \(a\) and \(h_0\)
replaced by \(b=(\log(1/a))^{-4}\): the KMT error is still \(o(1)\) after
normalization, the positive-half-line event \(\{V\le K\}\) has the same fixed
probability, and the level shift in
Lemma~\ref{lem:sech-persistence} remains \(o(1)\). Thus
\[
    \mathbb P\left\{
        S\left(\frac12+h\right)\le0
        \text{ for every }a\le h\le b
    \right\}
    \ge
    \exp\!\left\{-O\!\left(1+\log\frac ba\right)\right\},
    \qquad b=(\log(1/a))^{-4},
\]
successively on the intervals $[h_j,h_{j+1}]$. Using FKG again gives
\[
    \mathbb P\left\{
        S\left(\frac12+h\right)\le0
        \text{ for every }h_0\le h\le h_*
    \right\}
    \ge
    \exp\{-O(\log(1/h_0))\}.
\]
Since $\log(1/h_0)=O(\log\log(1/\beta))=o(\log(1/\beta))$, this factor is $\exp\{-o(\log(1/\beta))\}$.

Finally,
$\mathbb P\{S(1/2+h)\le0\text{ for every }h_*\le h<1/2\}>0$ by the
positive-probability construction used for $p_-(\alpha,\eta)>0$ in
Theorem~\ref{thm:polynomial-stable-regime}, with $b=1/2+h_*$. Multiplying
the three lower bounds gives $F(\beta)\ge\beta^{3/16+o(1)}$. Together with
the matching upper bound, this proves \eqref{eq:F-beta-persistence}; the
identities $q(\alpha)=F(\beta)$ and $p(\alpha)=F(\beta)^2$ finish the proof.
\end{proof}

Proposition~\ref{prop:upper-endpoint-persistence} implies
\(p(r)\to0=p(1/4)\) as \(r\uparrow1/4\). Together with the interior
continuity proved in Section~\ref{sec:phase-diagram}, this completes the proof
of Proposition~\ref{prop:phase-curve-properties}.
